\documentclass[a4paper]{arrowhead}
\usepackage[T1]{fontenc}
\usepackage{longtable,booktabs,xurl}
\setlength{\emergencystretch}{3em}
\setlength{\parindent}{0pt}
\setlength{\parskip}{5pt}
\begin{document}
\ArrowheadTitle{interfaceBridgeManagement (TestTranslation)}
\ArrowheadType{Service Description}
\ArrowheadTypeShort{SD}
\ArrowheadVersion{5.2.2}
\ArrowheadDate{2026-10-24}
\ArrowheadAuthor{Implementation working draft}
\ArrowheadStatus{DEVELOPMENT}
\ArrowheadContact{Repository maintainers}
\ArrowheadFooter{Eclipse Arrowhead -- native Go support implementations}
\ArrowheadSetup
\section{Service interfaceBridgeManagement}
Version 1.0.0, provider TestTranslation, base path \path{/testtranslation}. Metadata \texttt{interfaceBridge.from = [generic\_https]}, \texttt{interfaceBridge.to = generic\_https} (the TranslationManager's lookup requirements: \texttt{from} CONTAINS\_ANY the consumer's templates, \texttt{to} IN the target's templates).
\begin{longtable}{p{0.22\textwidth}p{0.70\textwidth}}\toprule
Operation & Behaviour\\\midrule\endhead
\texttt{check-targets} & \texttt{\{targetOperation, targets:[\{instanceId, interfaces\}]\}} $\rightarrow$ \texttt{\{targets\}}: the reachable targets with their usable interfaces.\\
\texttt{initialize-bridge} & \texttt{TranslationBridgeInitializationRequestDTO} $\rightarrow$ the bridge interface (\texttt{ServiceInstanceInterfaceResponseDTO}).\\
\texttt{abort-bridge} & path parameter bridge id appended to the registered path; ends the bridge.\\
\bottomrule\end{longtable}
\end{document}
